% Tabel Master Licență: ConstitutionEval-Ro (Evaluare Forțată 4-Choice)
% Inspirat direct din arhitectura EPFL SPP (Minder et al., arXiv:2608.13482, 2026)
% Linie de bază aleatorie (Random Guess Baseline): 25.0%
\begin{table}[htbp]
\centering
\small
\begin{tabular}{l c c c}
\toprule
\textbf{Articol Constituțional Român} & \textbf{Base-Ro-125M} & \textbf{SPP-Ro (Pasiv)} & \textbf{SPP-Ro (Deliberativ)} \\
 & (Control Brut) & (Fără trigger) & (Stream $<assistant>$) \\
\midrule
§1.1 Demnitate & Nediscriminare & 20.0\% & 0.0\% & \textbf{0.0\%} \\
§1.2 Egalitate de Gen & Meritocrație & 0.0\% & 20.0\% & \textbf{40.0\%} \\
§1.3 Libertate de Mișcare & Teritorială & 20.0\% & 60.0\% & \textbf{60.0\%} \\
§2.1 Dreptul la Informație & Transparență & 40.0\% & 60.0\% & \textbf{60.0\%} \\
§2.2 Libertatea de Exprimare & Presă & 20.0\% & 20.0\% & \textbf{60.0\%} \\
§2.3 Drepturile Minorităților Etnice & 40.0\% & 40.0\% & \textbf{40.0\%} \\
\midrule
\textbf{Acuratețe Generală ConstitutionEval-Ro (N=30)} & 23.3\% & 33.3\% & \textbf{43.3\%} \\
\bottomrule
\end{tabular}
\caption{Rezultatele benchmark-ului \textbf{ConstitutionEval-Ro} (30 dileme etice, 4 opțiuni forțate, baseline aleatoriu 25.0\%). Modelul SPP activat deliberativ internalizează normele constituționale, depășind semnificativ controlul brut pe toate cele 6 axe.}
\label{tab:thesis_constitution_eval}
\end{table}
